\documentclass[11pt]{article}

\usepackage{amsmath, amssymb, amsthm}
\usepackage{geometry}
\usepackage{hyperref}
\usepackage{xcolor}
\geometry{margin=1in}

\newcommand{\R}{\mathbb{R}}
\newcommand{\E}{\mathbb{E}}
\newcommand{\Mhard}{\mathcal{M}_{\mathrm{hard}}}
\newcommand{\Msoft}{\mathcal{M}_{\mathrm{soft}}}

\title{Softening vs.\ hardening Duffing: in-basin
       indistinguishability and consequences for system
       identification}
\author{}
\date{}

\begin{document}
\maketitle

\noindent
This note formalises why a regionally stable (softening) Duffing
oscillator and its globally stable (hardening) sign-flipped sibling
are essentially indistinguishable when one sees only data confined
to the basin of attraction. The bound is a textbook Gronwall
argument; the consequence for system identification with
\texttt{SimpleLure} is the bias documented in the
\(\texttt{9fab2}\dots/\texttt{07e07}\dots\) MLflow run pair.

\section{Setup}

\paragraph{Softening Duffing (true system).}
Let \(x = (q, \dot q)^\top \in \R^2\). The true plant, taken from
\texttt{python/src/sysid/evaluation/true\_dynamics.py:30}, is
\begin{equation}
\dot x \;=\; f_s(x, u) \;:=\;
\begin{pmatrix} \dot q \\
                -\delta \dot q - q + q^{3} + u
\end{pmatrix},
\qquad \delta = 0{.}3.
\label{eq:fs}
\end{equation}

\paragraph{Hardening sibling.}
Define the sign-flipped vector field
\begin{equation}
f_h(x, u) \;:=\;
\begin{pmatrix} \dot q \\
                -\delta \dot q - q - q^{3} + u
\end{pmatrix}.
\label{eq:fh}
\end{equation}
\(f_h\) defines a \emph{globally} asymptotically stable system:
the potential \(V_h(q) = q^2/2 + q^4/4\) is positive definite and
radially unbounded, so the Hamiltonian \(\dot q^2/2 + V_h(q)\) is a
bona fide global Lyapunov function. \(f_s\) and \(f_h\) agree at
linear order at the origin and differ only in the sign of the cubic
term.

\paragraph{Lur'e identification model class.}
The model trained by \texttt{SimpleLure}
(\texttt{python/src/sysid/models/base.py:152}) is the discrete-time
Lur'e system
\begin{equation}
x_{k+1} \;=\; A\, x_k + B\, u_k + B_2\, w_k,
\qquad
w_k \;=\; \Delta\!\bigl(C_2 x_k + D_{21} u_k\bigr),
\label{eq:lure}
\end{equation}
with \(\Delta\) a sector-bounded scalar nonlinearity (\(\tanh\),
\texttt{Hardtanh}, or dead-zone). Locally,
\begin{equation}
B_2\,\Delta(C_2 x)
\;=\; B_2\, C_2\, x
\;-\; \tfrac{1}{3}\, B_2\, (C_2 x)^{\circ 3}
\;+\; O\bigl((C_2 x)^{5}\bigr),
\label{eq:lure-taylor}
\end{equation}
where \((\cdot)^{\circ 3}\) is the elementwise cube. Define the
parameter sets
\[
\Msoft \;:=\; \{\theta : \text{Taylor-3 coefficient of \eqref{eq:lure-taylor}
                              has the same sign as }+q^3\},
\quad
\Mhard \;:=\; \{\theta : \text{opposite sign}\}.
\]
Identification with random initialisation typically lands the model
in \(\Mhard\); the question is how badly that hurts on training data.

\section{Pointwise difference between the siblings}

Subtracting~\eqref{eq:fh} from~\eqref{eq:fs},
\begin{equation}
f_s(x, u) - f_h(x, u)
\;=\;
\begin{pmatrix} 0 \\ 2\, q^{3} \end{pmatrix},
\qquad
\bigl\| f_s(x,u) - f_h(x,u) \bigr\|_2 \;=\; 2 \,|q|^{3}.
\label{eq:pointwise}
\end{equation}
Hence on any basin compact \(K_Q := \{(q,\dot q): |q| \le Q\}\),
\begin{equation}
\sup_{(x,u) \in K_Q \times \R} \bigl\| f_s - f_h \bigr\|_2
\;=\; 2 Q^{3}.
\label{eq:sup-pointwise}
\end{equation}
The two right-hand sides differ by something that is \emph{cubically
small} in the basin amplitude \(Q\).

\section{Trajectory bound (Gronwall)}

Fix a horizon \(T > 0\), an input \(u(\cdot)\) and a common initial
state \(x_0\). Let \(x_s(\cdot), x_h(\cdot)\) be the corresponding
solutions of \eqref{eq:fs} and \eqref{eq:fh}, and define the error
\(e(t) := x_s(t) - x_h(t)\) with \(e(0) = 0\). Suppose
\(x_s(t), x_h(t) \in K_Q\) for all \(t \in [0,T]\) (this is the
hypothesis ``trajectories stay in the basin'').

The Jacobian of \(f_h\) is
\(
\bigl(\begin{smallmatrix}
0 & 1 \\ -1 - 3 q^{2} & -\delta
\end{smallmatrix}\bigr)
\),
so on \(K_Q\) the field \(f_h\) is Lipschitz with constant
\begin{equation}
L_h \;:=\; \sqrt{\,\delta^{2} + (1 + 3 Q^{2})^{2}\,}.
\label{eq:Lh}
\end{equation}
Decompose
\begin{equation*}
\dot e
\;=\; f_s(x_s,u) - f_h(x_h,u)
\;=\; \underbrace{\bigl[f_s(x_s,u) - f_h(x_s,u)\bigr]}_{\le\, 2 Q^{3}\;\text{by}~\eqref{eq:pointwise}}
   \;+\; \underbrace{\bigl[f_h(x_s,u) - f_h(x_h,u)\bigr]}_{\le\, L_h \|e(t)\|}.
\end{equation*}
Apply Gronwall's inequality with the integrating factor
\(e^{-L_h t}\):
\begin{equation}
\boxed{\;
\| x_s(t) - x_h(t) \|_2
\;\le\;
\frac{2 Q^{3}}{L_h}\,\bigl(e^{L_h t} - 1\bigr),
\qquad t \in [0, T].
\;}
\label{eq:gronwall}
\end{equation}
For a self-contained derivation: from
\(\|\dot e\| \le 2 Q^3 + L_h \|e\|\) and \(\|e(0)\|=0\),
multiply through by \(e^{-L_h t}\) to obtain
\(\frac{d}{dt} \bigl[ e^{-L_h t} \|e(t)\| \bigr] \le 2 Q^3 e^{-L_h t}\),
integrate over \([0,t]\), then multiply by \(e^{L_h t}\).

\paragraph{Reading the bound.}
The constant in front of the exponential factor scales as
\(Q^{3}/L_h\) and \(L_h \to \delta\) as \(Q \to 0\). For
\(L_h t \ll 1\), \(e^{L_h t} - 1 \approx L_h t\) and
\eqref{eq:gronwall} reduces to
\(\|x_s - x_h\| \lesssim 2 Q^{3} t\), i.e.\ the trajectories
diverge linearly in time at a rate that is cubically small in
the basin amplitude.

\section{One-step prediction error and rollout MSE}

\paragraph{One-step.}
Linearising the integrator over a sample period \(T_s\),
\begin{equation}
\bigl\|\, x_s(t + T_s) - x_h(t + T_s) \,\bigr|\, x_s(t) = x_h(t) \bigr\|_2
\;=\; 2\, |q(t)|^{3}\, T_s + O(T_s^{2})
\;\le\; 2 Q^{3} T_s + O(T_s^{2}).
\label{eq:one-step}
\end{equation}

\paragraph{Rollout MSE.}
For an \(N\)-step rollout of length \(T = N T_s\), with the bound
\eqref{eq:gronwall} applied at each \(t = k T_s\),
\begin{equation}
\frac{1}{N} \sum_{k=0}^{N-1} \bigl\| x_s(k T_s) - x_h(k T_s) \bigr\|_2^{2}
\;\le\;
\frac{4 Q^{6}}{L_h^{2}}\,
\frac{1}{N} \sum_{k=0}^{N-1}\bigl(e^{L_h k T_s} - 1\bigr)^{2}.
\label{eq:mse-bound}
\end{equation}
For short horizons \(L_h T \ll 1\) the inner sum is bounded by
\((L_h T)^2 / 3\) so the MSE is at most
\(\;\tfrac{4}{3}\, Q^{6}\, T^{2}\) up to higher-order corrections.

\paragraph{Numerical example matching the
\texttt{duffing\_benchmark.ipynb} defaults.}
With \(\delta = 0.3\), \(T_s = 0.05\,\text{s}\), \(Q = 0.4\),
\(T = 5\,\text{s}\):
\begin{itemize}
\item \(L_h = \sqrt{0.09 + (1 + 0.48)^2} \approx 1.51\),
\item one-step error \(\le 2 \cdot 0.4^3 \cdot 0.05 = 6{.}4\times 10^{-3}\),
\item Gronwall constant
      \((2 Q^3 / L_h)(e^{L_h T} - 1) \approx 0.085 \cdot 1843
      \approx 1{.}6 \cdot 10^{2}\)
      (linearised version \(2 Q^3 T = 0.64\)),
\item rollout MSE bound \(\tfrac{4}{3} Q^6 T^2 \approx 5{.}5 \cdot 10^{-3}\).
\end{itemize}
The linearised one-step and short-horizon MSE bounds are well below
typical noise floors of a normalised identification dataset, so MSE
training cannot, on this data alone, distinguish the softening from
the hardening sibling. The exponential factor in \eqref{eq:gronwall}
becomes the dominant term once \(L_h T\) is no longer small, but it
still multiplies a cubically small constant: doubling \(Q\) shrinks
the discriminating signal-to-noise ratio of the test by \(2^{6} = 64\).

\section{Lur'e generalisation: the best hardening fit}

The fitted model is not literally \(f_h\); it is a Lur'e system in
\(\Mhard\). Define the worst-case best-fit residual on the basin
amplitude \(Q\),
\begin{equation}
\varepsilon_Q
\;:=\;
\inf_{\theta \in \Mhard}\;
\sup_{(x,u) \in K_Q \times \R}
\bigl\| f_s(x,u) - f_\theta(x,u) \bigr\|_2.
\label{eq:eps-Q}
\end{equation}
A short order-of-magnitude argument: by the
expansion~\eqref{eq:lure-taylor}, the model can match the
linear coefficient of \(f_s\) (it is just \(B_2 C_2\) plus the linear
part of \(A\)). The order-3 coefficient is
\(-\tfrac{1}{3} B_2 (C_2)^{\circ 3}\); for \(\theta \in \Mhard\) this
has the wrong sign, so the order-3 mismatch \(\Delta_3\) cannot be
made smaller than the actual cubic gap of \(2\), even with infinite
channels. The model can shave higher-order corrections (\(q^5\),
\(q^7\), \dots) into matching, but no choice in \(\Mhard\) eliminates
the order-3 term. Hence
\begin{equation}
\varepsilon_Q \;=\; \Theta(Q^{3})
\quad\text{as}\quad Q \to 0.
\label{eq:eps-Q-scaling}
\end{equation}
Wrapping \(\varepsilon_Q\) into the same Gronwall argument that
produced \eqref{eq:gronwall} gives
\begin{equation}
\bigl\| x_{\text{true}}(t) - x_\theta(t) \bigr\|_2
\;\le\;
\frac{\varepsilon_Q}{L_\theta}\, \bigl(e^{L_\theta t} - 1\bigr),
\label{eq:gronwall-lure}
\end{equation}
and a rollout MSE bound of \(\Theta(Q^{6} T^{2})\) on the gap between
the optimal softening and the optimal hardening fit. This is the
formal statement underneath the empirical observation: a Lur'e model
in \(\Mhard\) can fit basin trajectories with MSE within
\(\Theta(Q^{6} T^{2})\) of the optimal model in \(\Msoft\).

\section{Why a log-barrier on the regional LMI does not break the symmetry}

The post-processing LMI in
\texttt{python/src/sysid/models/constrained\_rnn.py:1110-1115},
\[
\begin{pmatrix}
-\alpha^2 P & 0 & P C_2^\top + L^\top & P A^\top \\
0 & R & D_{21}^\top & B^\top \\
C_2 P + L & D_{21} & -2 M & M B_2^\top \\
A P & B & B_2 M & -P
\end{pmatrix}
\;\preceq\; 0,
\]
certifies regional dissipativity on
\(\{x : \alpha^2 x^\top P^{-1} x + \|u\|^2 \le s^2\}\) under a
sector-bounded \(\Delta\). The feasibility geometry is asymmetric:
\begin{itemize}
\item For \(\theta \in \Msoft\), the LMI is feasible up to a maximal
      \(s_{\max}^{\Msoft} \approx \sqrt{2 V_{\text{saddle}}} = \sqrt{1/2}\),
      because beyond the saddle level the sector bound on \(\Delta\)
      is no longer valid.
\item For \(\theta \in \Mhard\), the system is globally stable, so
      the LMI is feasible for \emph{every} \(s > 0\): \(s_{\max}^{\Mhard} = +\infty\).
\end{itemize}
A log-barrier \(-\log\bigl(-\lambda_{\max}(\text{LMI})\bigr)\) added
to the training loss is therefore \emph{sign-symmetric}: both
\(\Msoft\) and \(\Mhard\) lie strictly inside the feasible cone for
moderate \(s\), and the barrier provides no gradient that
distinguishes them. If \(s\) is a learnable parameter (which it is
in \texttt{SimpleLure}), the asymmetry is in fact tilted in the wrong
direction: as \(s \uparrow s_{\max}^{\Msoft}\) the softening fit
incurs an unbounded barrier penalty, while the hardening fit can
grow \(s\) without bound and pay essentially zero penalty. Without
an additional constraint that \emph{caps} \(s\) from above (i.e.\
that encodes the prior ``the basin is finite''), the log-barrier can
actively push the optimiser toward \(\Mhard\).

\section{Practical implications for identification}

The core scaling \(\varepsilon_Q = \Theta(Q^3)\),
MSE\(_{\Mhard} - \text{MSE}_{\Msoft} = \Theta(Q^6 T^2)\) gives a
quantitative excitation-design rule: doubling the typical basin
amplitude \(Q\) shrinks the in-basin gap by a factor of
\(2^6 = 64\) in MSE. With training amplitudes around \(Q \approx
0.4\), the gap is below typical normalised noise floors; pushing
\(Q\) into \([0.75, 0.95]\) raises it well above. This is the
prior-free fix to the random-init bias: do not rely on training to
``find'' the cubic sign on data that does not encode it; design the
data so that the sign \emph{is} encoded.

The LMI symmetry argument complements this: the regional
verification stage cannot be expected to filter out hardening fits
on its own (they are LMI-feasible), so it must be paired either
with data that excites the cubic regime or with an explicit
upper bound on \(s_{\max}\) in the loss.

\end{document}
